\clearpage
\appendix
\counterwithin{equation}{section}
\renewcommand{\theequation}{\thesection.\arabic{equation}}

\section{Proofs}
\label{app:proofs}

\subsection{Count law and conditional information}
\label{app:model}
\label{app:target1}

The probability generating function of one negative binomial offspring count is $[1+(R_t/k)(1-z)]^{-k}$. Conditional independence raises this function to the power $n$. Substituting $z\mapsto1-\rho+\rho z$ for independent reporting gives $[1+(\rho R_t/k)(1-z)]^{-nk}$, proving Eq.~(\ref{eq:observed}). Its derivatives give the stated moments.

For fixed, known $(n,\rho,k)$ and $R>0$, the count likelihood has score and Fisher information
\begin{equation}
 \dot\ell(R)=\frac{c}{R}-\frac{(c+nk)\rho}{k+\rho R}
 =\mathcal I(R)(\Rhat-R),\qquad
 \mathcal I(R)=\frac{m}{R+\rho R^2/k}.
 \label{eq:score}
\end{equation}
Consequently $\Var(\Rhat\mid R)=\mathcal I(R)^{-1}$. The support does not depend on $R$, and the regular one-parameter Cram\'er--Rao inequality applies to estimators unbiased for every $R$. It gives the attained bound in Section~\ref{sec:background} \citep{lehmann1998theory}. The calculation does not establish an unrestricted decision-theoretic bound or a bound with unknown nuisance parameters.

\subsection{Proof of Proposition~\ref{prop:average}}
\label{app:prop1}

Let $e_t=\Rhat_t-R_t$. Then $\E(e_t\mid R_t)=0$ and
\begin{equation}
 \E(e_t^2)=\frac{\E R_t+\rho\E R_t^2/k}{m}=\frac{A}{m},
 \qquad \E[e_t(R_t-\Rbar)]=0.
\end{equation}
Squaring Eq.~(\ref{eq:error-split}) proves Eq.~(\ref{eq:var-average}). For any estimator based on the same count and unbiased conditional on every $R_t$, its expected conditional variance is at least $\E[\mathcal I(R_t)^{-1}]=A/m$. Its process-mean risk adds the same $v_R$, proving the restricted optimality statement. At fixed $\rho$, $k$ and transmission law,
\begin{equation}
 \frac{V(m)-V(2m)}{V(m)}=\frac{1}{2(1+m/\mx)},\qquad v_R>0.
\end{equation}
This comparison changes exposure $n$; it does not hold unchanged when reporting probability changes.

\subsection{Proof of Corollary~\ref{cor:targets}}
\label{app:cor-targets}

Conditional on the entire path $R_{1:T}$, the $e_t$ are independent and centred. They are therefore pairwise uncorrelated and orthogonal to every square-integrable function of that path. Thus
\begin{equation}
 \E\!\left[\left(T^{-1}\sum_t e_t\right)^2\right]=\frac{A}{Tm},\qquad
 \Var(\widehat{\Rbar}_T)=\Var(\Rbar_T)+\frac{A}{Tm}.
\end{equation}
For independent, identically distributed $R_t$, $\Var(\Rbar_T)=v_R/T$. Under the stated stationary autocovariance, summing all covariance pairs gives
$T^2\Var(\Rbar_T)=v_R[T+2\sum_{h=1}^{T-1}(T-h)\phi^h]$.
The finite geometric sum yields Eq.~(\ref{eq:ar-sum}), including $T=1$ and $\phi=0$ \citep{hamilton1994time}. Similarly, summing spatial covariance pairs gives Eq.~(\ref{eq:pooling}); its correlation matrix must be admissible.

\subsection{Proof of Proposition~\ref{prop:report-noise} and Corollary~\ref{cor:report-current}}
\label{app:prop2}
\label{app:cor-report-current}

Write $z_t=p_t/\rho$ and $u_t=C_t/m-z_tR_t$. Conditional on $(R_t,p_t)$, $u_t$ has zero mean. Independence of $p_t$ and $R_t$, with $\E z_t=1$ and $\E z_t^2=1+\tau^2$, gives
\begin{equation}
 \E(u_t^2)=\frac{\Rbar+\rho(1+\tau^2)(\Rbar^2+v_R)/k}{m}
 =\frac{A_H}{m}.
\end{equation}
The error $\Rhat_t-R_t=u_t+(z_t-1)R_t$ has uncorrelated terms. The second has expected square $\tau^2(\Rbar^2+v_R)=v_H$, proving Eq.~(\ref{eq:vH-current}). Moreover $\E(\Rhat_t-R_t\mid R_t)=0$, so process-mean risk adds $v_R$. Under the additional cross-period assumptions of Corollary~\ref{cor:report-current}, these errors are centred and independent conditional on the transmission path after integrating out reporting, giving Eq.~(\ref{eq:report-means}).

For the identification claim, set $U=(1+\tau^2)(\Rbar^2+v_R)$. The two observable moments of the scaled count satisfy
\begin{equation}
 \E\Rhat_t=\Rbar,\qquad
 \Var(\Rhat_t)=U-\Rbar^2+\frac{\Rbar+\rho U/k}{m}.
 \label{eq:moment-identification}
\end{equation}
Whenever admissible distinct values of $\tau^2$ give $v_R=U/(1+\tau^2)-\Rbar^2\ge0$, these moments coincide. Such distinct pairs occur away from degenerate boundaries; varying exogenous $m$ does not separate them. The restrictions on $p_t$ and $R_t$ must still hold. This proves nonidentification by these moments, without asserting equality of full count distributions.

For $v_R>0$, the generalised balance scale $m_{\times,H}=A_H/(v_R+v_H)$ obeys
\begin{equation}
 \mx-m_{\times,H}
 =\frac{\tau^2(\Rbar^2+v_R)\Rbar(1+\rho\Rbar/k)}{v_R(v_R+v_H)}\ge0.
 \label{eq:mx-gap}
\end{equation}
For $\Rbar>0$, equality holds exactly when $\tau=0$. A lower balance scale here accompanies a larger risk, not more accurate surveillance.

\section{Symbolic theoretical extensions}
\label{app:extensions}

\subsection{Forecasting and shrinkage}
\label{app:forecast}
\label{app:shrinkage}

Under fixed, known reporting and the stationary covariance model of Corollary~\ref{cor:targets}, assume the counting error at $t$ is centred conditional on the entire transmission path. Then
\begin{equation}
 \E[(\Rhat_t-R_{t+1})^2]=\frac{A}{m}+2(1-\phi)v_R.
\end{equation}
For known $\Rbar$, $v_R$ and $\phi$, minimising over forecasts of the form $\Rbar+\lambda(\Rhat_t-\Rbar)$ gives
\begin{equation}
 \lambda^*=\frac{\phi v_R}{v_R+A/m},\qquad
 \mathrm{MSE}_{\mathrm{lin}}=v_R-\frac{\phi^2v_R^2}{v_R+A/m}.
 \label{eq:forecast}
\end{equation}
These follow by expanding a quadratic in $\lambda$ and using $\Cov(\Rhat_t,R_{t+1})=\phi v_R$. The limit is $(1-\phi^2)v_R$; with independent transmission periods, the best forecast in this class is the known process mean. Optimality is restricted to linear forecasts from the single current estimate, not all histories or nonlinear predictors.

For estimating $\Rbar$ instead, consider $\hat\theta_\lambda=(1-\lambda)\Rhat_t+\lambda\mu_0$ with fixed benchmark $\mu_0$ and $\Delta=\mu_0-\Rbar$. Its risk is
\begin{equation}
 \E[(\hat\theta_\lambda-\Rbar)^2]=(1-\lambda)^2V+\lambda^2\Delta^2,
 \qquad V=v_R+A/m.
\end{equation}
The oracle weight $\lambda^*=V/(V+\Delta^2)$ gives risk $V\Delta^2/(V+\Delta^2)$. This may be below $v_R$, illustrating the importance of the unbiasedness restriction. The oracle uses unknown quantities; estimating its weight requires a separate risk analysis.

\subsection{Budget allocation with feasible periods}
\label{app:allocation}

\begin{proposition}[Continuous budget relaxation]
\label{prop:allocation}
\label{app:prop3}
Under H1--H4, independent, identically distributed transmission periods, conditionally independent counts, fixed $\rho$ and linear costs, let $\mathcal B=T(c_0+c_1m)$ with $c_0,c_1>0$. Treat $T\ge1$ and $m>0$ as continuous. If $v_R>0$ and $\mathcal B>c_0$, the unconstrained stationary point is
\begin{equation}
 m^*=\sqrt{\frac{c_0A}{c_1v_R}}=\sqrt{\frac{c_0}{c_1}\mx}.
 \label{eq:optimal-m}
\end{equation}
The feasible optimum is $\min\{m^*,(\mathcal B-c_0)/c_1\}$, with $T$ recovered from the budget constraint. A prescribed lower bound $m_{\min}>0$ instead requires clamping $m^*$ to $[m_{\min},(\mathcal B-c_0)/c_1]$, provided this interval is nonempty.
\end{proposition}

\begin{proof}
For independent, identically distributed periods, process-mean risk is
\begin{equation}
 \frac{v_R}{T}+\frac{A}{Tm}
 =\frac{(c_0+c_1m)(v_R+A/m)}{\mathcal B}=:f(m).
 \label{eq:alloc-indep}
\end{equation}
Its derivatives are $f'(m)=(c_1v_R-c_0A/m^2)/\mathcal B$ and $f''(m)=2c_0A/(\mathcal B m^3)>0$. Strict convexity gives the stationary point and the clamped solution. When $v_R=0$, risk decreases with $m$, so the optimum is the largest feasible $m$, corresponding to $T=1$ without further constraints.
\end{proof}

The relaxation has the familiar form of cluster-size allocation \citep{cochran1977sampling, raudenbush1997statistical}. Actual designs require integer $T$ and integer $n$, with $m=\rho n$. Under a budget inequality, for each feasible integer $T$ choose
\begin{equation}
 n_T=\left\lfloor\frac{\mathcal B/T-c_0}{c_1\rho}\right\rfloor\ge1,
 \qquad m_T=\rho n_T,
\end{equation}
and compare the risks over all feasible $T$. This handles unspent budget and the boundary $T=1$; rounding the relaxed optimum alone need not suffice.

\subsection{Budget allocation under temporal dependence}
\label{app:budget-ar}

For the stationary autocovariance model, compare instead the exact integer-design objective
\begin{equation}
 F(T)=\frac{v_RS(\phi,T)}{T^2}+\frac{A}{Tm_T}.
 \label{eq:budget-ar-exact}
\end{equation}
No independent-period optimum or universal switch threshold follows from this expression. Under independent reporting shocks, replace its counting term by $(v_H+A_H/m_T)/T$. All costs and exposures must refer to the same period; these symbolic formulas provide no direct conversion of weekly admission ratios into generational design thresholds.

\subsection{Coverage and reporting completeness}
\label{app:coverage}

For fixed transmission law and a differentiable reporting-volatility relation $\tau^2(\rho)$ independent of $n$, write $Q=\Rbar^2+v_R$. At interior values of $\rho$, the reductions in single-period process-mean risk from proportional increases in $n$ and $\rho$ are
\begin{equation}
\begin{aligned}
 -\frac{\partial V_H}{\partial\log n}
 &=\frac{\Rbar}{\rho n}+\frac{(1+\tau^2)Q}{kn},\\
 -\frac{\partial V_H}{\partial\log\rho}
 &=\frac{\Rbar}{\rho n}-\rho\frac{\mathrm d\tau^2}{\mathrm d\rho}\,Q\left(1+\frac{1}{kn}\right).
\end{aligned}
 \label{eq:coverage-gains}
\end{equation}
These follow by differentiating Eq.~(\ref{eq:vH-AH}) before substituting $m=\rho n$. At fixed $\tau$, the first gain is larger. If improving completeness also stabilises reporting, that ordering can change; these are proportional, not cost-adjusted, comparisons.

\begin{corollary}[A specified reporting mechanism]
\label{cor:switch}
If $p_t$ is the fraction of $\nu$ independent reporting units participating with probability $\rho$, so $\tau^2=(1-\rho)/(\rho\nu)$, increasing completeness has the larger proportional gain exactly when
\begin{equation}
 m>\frac{\rho^2(\nu-1)}{k}.
 \label{eq:r3-switch}
\end{equation}
For $p_t\sim\mathrm{Beta}(\rho\nu,(1-\rho)\nu)$, the corresponding condition is $m>\rho^2\nu/k$.
\end{corollary}

Subtracting the first derivative gain from the second gives $Q(kn)^{-1}[(kn+\rho)/(\rho\nu)-1]$ for the reporting-unit model. Replacing $\nu$ by $\nu+1$ proves the beta case. Neither mechanism is implied by observed facility participation alone.

\subsection{Conditional renewal moments and noisy denominators}
\label{app:renewal}
\label{app:random-denominator}

A restricted renewal analogue can be stated explicitly. Given history $\mathcal F$, suppose infection pressure $\Lambda>0$ is independent of the next $R$, with $\E(R\mid\mathcal F)=\Rbar$ and $\Var(R\mid\mathcal F)=v_R$. Assume a conditional count law with mean $\rho\Lambda R$ and variance $\rho\Lambda R+\rho^2\Lambda R^2/k$. Put $\bar\Lambda=\E(\Lambda\mid\mathcal F)$ and $m=\rho\bar\Lambda$. Total variance gives
\begin{equation}
 \frac{\Var(C\mid\mathcal F)}{m^2}
 =v_R+\frac{A}{m}
   +\frac{\Var(\Lambda\mid\mathcal F)}{\bar\Lambda^2}(\Rbar^2+v_R).
 \label{eq:renewal-conditional}
\end{equation}
The result assumes the conditional independence and fixed moments just stated. It is a count-prediction variance, not a general error identity for a ratio with an estimated denominator, and it does not resolve endogenous incidence.

To see the ratio issue, condition on a positive latent path and write numerator and denominator as their latent values times $(1+\varepsilon)$ and $(1+u)$. Let $Y=y(1+\varepsilon)/(1+u)$, with centred disturbances of order $\tau$, a denominator bounded away from zero, bounded scaled third moments and integrable Taylor remainders. Then
\begin{equation}
\begin{aligned}
 \E(Y\mid y,\text{path})/y
 &=1+\E(u^2\mid\text{path})-\E(\varepsilon u\mid\text{path})+O(\tau^3),\\
 \Var(Y\mid\text{path})/y^2
 &=\Var(\varepsilon-u\mid\text{path})+O(\tau^3).
\end{aligned}
\end{equation}
Both numerator--denominator covariance and induced bias enter apparent variation. Overlapping ratios can share disturbances. These conditional approximations do not supply an exact source decomposition of a fitted loss intercept.

\subsection{Replication and variance estimation}
\label{app:vr-precision}
\label{app:seasons}

\begin{proposition}[A conditional information bound for process variance]
\label{prop:vr-precision}
Suppose $R_1,\ldots,R_T$ are independent, identically distributed gamma variables with known mean $\Rbar$ and variance $v_R>0$. If the observation law given $R_{1:T}$ does not otherwise depend on $v_R$, any regular unbiased estimator of $v_R$ from the counts has relative standard error at least
\begin{equation}
 [T\xi(\xi\psi'(\xi)-1)]^{-1/2},\qquad \xi=\Rbar^2/v_R,
 \label{eq:vr-bound}
\end{equation}
where $\psi'$ is the trigamma function.
\end{proposition}

Observing the latent $R_t$ cannot provide less information than observing only counts generated from them. With the mean fixed, information for gamma shape $\xi$ is $\psi'(\xi)-1/\xi$ per period; the change of parameter $\mathrm d\xi/\mathrm dv_R=-\xi/v_R$ gives the bound. It approaches $\sqrt{2/T}$ as $\xi\to\infty$.

\begin{corollary}[Independent season effects]
\label{cor:season}
For $S\ge2$ independent, identically distributed Gaussian season effects with unknown mean and variance $v_S>0$, assume their conditional count law is parameter-free. Any estimator unbiased for $v_S$ for all means and variances has relative standard error at least $\sqrt{2/(S-1)}$.
\end{corollary}

Conditional expectation onto the latent effects cannot increase variance. Their sample variance is uniformly minimum-variance unbiased and has variance $2v_S^2/(S-1)$ \citep{lehmann1998theory, searle1992variance}. Thus more observations within existing seasons cannot replace independent seasons for this bound. Both bounds are conditional model results, not empirical confidence limits.

\FloatBarrier
\clearpage
\renewcommand{\thetable}{S\arabic{table}}
\setcounter{table}{0}
\renewcommand{\thefigure}{S\arabic{figure}}
\setcounter{figure}{0}
\section{Supplementary methods and results}
\label{app:supplement}

\subsection{Input registry and reproducibility}
\label{app:inputs}

All reported analyses are regenerated from the raw NHSN state-week table whose SHA-256 is recorded in the release manifest. The panel reconstructs missing values as missing and requires four consecutive weekly dates; it does not insert zero admissions. The primary samples and all fixed windows are identified by deterministic keys in the release. Every retained table, figure and inline macro is generated from the same versioned result bundle; the manifest records input hashes, generator hash, actual bootstrap budgets, fit failures and Monte Carlo diagnostics. A resumed computation must match the complete source, input and configuration fingerprint; otherwise it is run afresh.

Table~\ref{tab:pooled} reports pooled affine working fits and the three generation-interval transforms. The 3.6-day and fixed-coefficient-of-variation variants are stated as choices; the 3.2-day, SD 2.1-day variant reproduces the intrinsic interval estimate of \citet{chan2025estimating}. Table~\ref{tab:cv} reports fold-level held-out scores from training-only centres and fits. The season folds use the same data as the fit and are descriptive.

The calibration registry declares each input as a documented external estimate, a value derived from the observed panel, or an analysis assumption. Delay, offspring dispersion, admission probability and the weekly shock variance are assumptions in the simulator; their structural role is sensitivity, not measurement. In the observation model, binomial admission allocation is the baseline and $\varphi>1$ describes direct gamma--Poisson counts with conditional variance $\varphi\lambda$. The simulator is checked for conditional observation moments and for the direct-count construction; those tests verify implementation, not empirical adequacy.

\subsection{Conditional calibration}
\label{app:calibration}

Each declared scenario is simulated on the four fixed observed window lengths with the observed first-week state counts. Initial zero counts remain zero and are not raised; this can reduce the retained sample relative to the empirical panel, and the discrepancy is retained in the response diagnostics. Weekly reproduction shocks separate national and state components; persistence in reporting is a separate mechanism. Reproduction-number innovations are truncated below at a declared positive bound, and the truncation fraction is recorded.

For each empirical statistic, a weighted line relates the simulated mean intercept to the declared weekly variance grid. Paired empirical draws and paired simulated panels are combined to obtain conditional inversion draws. Failed fits, nonfinite values, nonpositive implied values and grid extrapolation are counted and reported; percentile intervals refer only to finite draws and are labelled conditional. Several scenarios can give similar empirical summaries while differing in observation mechanism. The calibration therefore reports response checks for intercept, slope, sample size, signed-deviation variance ratio and within-state and national lag-one autocorrelation. It does not assign a definitive biological variance or a probability that one mechanism is correct.

\subsection{Secondary sensitivities and design illustration}
\label{app:secondary}

Table~\ref{tab:secondary} reports season-specific intercepts and a paired participation adjustment in which each statistic's counts are divided by facility participation while the original scale is retained. The adjustment measures only the recorded participation proxy; testing intensity, coding and care seeking inside reporting facilities are not observed. The three summer summaries are unadjusted and describe a single voluntary period; they are not causal comparisons or independent replications. No equivalence or mandate effect is claimed.

Table~\ref{tab:design} reports a model-conditional illustration of curve discrimination. Observed season designs are recycled with declared count multipliers and crossed mean-one lognormal variation; the generating curves are the fitted affine and power-law means. Calibration quantiles and evaluation simulations are independent, and the null rejection fraction is reported with an exact binomial interval before power is interpreted. Power was near one in every cell. That saturation shows the synthetic generator is far more favourable than the observed data, where the two curves were not distinguishable; the exercise therefore does not answer how many seasons would be needed in practice and is reported only to document that limitation. The budget expressions in Appendix~\ref{app:allocation} remain symbolic; no weekly admission ratio is converted into a generational design threshold.

\begin{table}[htbp]
\centering
\small
\caption{Pooled affine loss fits and generation-interval sensitivities}\label{tab:pooled}
\setlength{\tabcolsep}{2pt}
\begin{tabularx}{\linewidth}{@{}>{\raggedright\arraybackslash}p{0.2400\linewidth}>{\raggedright\arraybackslash}p{0.1750\linewidth}>{\raggedright\arraybackslash}p{0.1750\linewidth}>{\raggedright\arraybackslash}p{0.1750\linewidth}>{\raggedright\arraybackslash}p{0.1750\linewidth}@{}}
\toprule
Statistic & Parameter & n & Estimate [95\% CI] & Finite/B \\
\midrule
RA & a & 1748 & 0.01495 [0.01001, 0.02092] & 2000/2000 \\
RA & b & 1748 & 0.56 [0.369, 0.7692] & 2000/2000 \\
RA & b/a & 1748 & 37.44 [18.93, 69.92] & 2000/2000 \\
G & a & 1756 & 0.08401 [0.04986, 0.1234] & 2000/2000 \\
G & b & 1756 & 1.971 [0.9889, 3.042] & 2000/2000 \\
G & b/a & 1756 & 23.46 [8.279, 54.7] & 2000/2000 \\
RA32 & a & 1748 & 0.01191 [0.007997, 0.01668] & 2000/2000 \\
RA32 & b & 1748 & 0.4482 [0.2975, 0.6147] & 2000/2000 \\
RA32 & b/a & 1748 & 37.62 [19.13, 69.89] & 2000/2000 \\
RA32source & a & 1748 & 0.01106 [0.007501, 0.01553] & 2000/2000 \\
RA32source & b & 1748 & 0.4196 [0.2827, 0.5688] & 2000/2000 \\
RA32source & b/a & 1748 & 37.93 [19.42, 69.26] & 2000/2000 \\
\bottomrule
\end{tabularx}
\par\smallskip
\begin{minipage}{\linewidth}\footnotesize\raggedright
Signed percentile intervals; ratios are defined only for positive denominators. Finite-draw counts are retained for every estimate. RA a: 2000/2000 finite draws; RA b: 2000/2000 finite draws; RA ratio: 2000/2000 finite draws; G a: 2000/2000 finite draws; G b: 2000/2000 finite draws; G ratio: 2000/2000 finite draws; RA32 a: 2000/2000 finite draws; RA32 b: 2000/2000 finite draws; RA32 ratio: 2000/2000 finite draws; RA32source a: 2000/2000 finite draws; RA32source b: 2000/2000 finite draws; RA32source ratio: 2000/2000 finite draws.
\end{minipage}
\end{table}

\begin{table}[htbp]
\centering
\scriptsize
\caption{Foldwise model validation}\label{tab:cv}
\setlength{\tabcolsep}{2pt}
\begin{tabularx}{\linewidth}{@{}>{\raggedright\arraybackslash}p{0.2400\linewidth}>{\raggedright\arraybackslash}p{0.1167\linewidth}>{\raggedright\arraybackslash}p{0.1167\linewidth}>{\raggedright\arraybackslash}p{0.1167\linewidth}>{\raggedright\arraybackslash}p{0.1167\linewidth}>{\raggedright\arraybackslash}p{0.1167\linewidth}>{\raggedright\arraybackslash}p{0.1167\linewidth}@{}}
\toprule
Statistic / split & Fold & Valid/total or n & Affine & Power & Difference & Wins / valid \\
\midrule
RA / state & Summary & 51/51 & -- & -- & -0.002018 & 26 \\
RA / season & Summary & 4/4 & -- & -- & 0.004626 & 3 \\
RA & 2022-23 & 365 & -2.378 & -2.369 & 0.008954 & yes \\
RA & 2023-24 & 499 & -2.613 & -2.612 & 0.0008806 & yes \\
RA & 2024-25 & 356 & -2.46 & -2.481 & -0.02077 & yes \\
RA & 2025-26 & 528 & -2.346 & -2.324 & 0.0223 & yes \\
G / state & Summary & 51/51 & -- & -- & -0.007369 & 22 \\
G / season & Summary & 4/4 & -- & -- & -0.004414 & 2 \\
G & 2022-23 & 367 & -0.5979 & -0.5942 & 0.003756 & yes \\
G & 2023-24 & 503 & -1.201 & -1.192 & 0.008487 & yes \\
G & 2024-25 & 357 & -0.7333 & -0.7584 & -0.02505 & yes \\
G & 2025-26 & 529 & -0.7557 & -0.7641 & -0.008421 & yes \\
RA32 / state & Summary & 51/51 & -- & -- & -0.002242 & 25 \\
RA32 / season & Summary & 4/4 & -- & -- & 0.004343 & 3 \\
RA32 & 2022-23 & 365 & -2.612 & -2.604 & 0.00805 & yes \\
RA32 & 2023-24 & 499 & -2.83 & -2.83 & 0.0001619 & yes \\
RA32 & 2024-25 & 356 & -2.691 & -2.712 & -0.02118 & yes \\
RA32 & 2025-26 & 528 & -2.572 & -2.549 & 0.02294 & yes \\
RA32source / state & Summary & 51/51 & -- & -- & -0.002824 & 25 \\
RA32source / season & Summary & 4/4 & -- & -- & 0.003577 & 2 \\
RA32source & 2022-23 & 365 & -2.701 & -2.695 & 0.005641 & yes \\
RA32source & 2023-24 & 499 & -2.886 & -2.888 & -0.001735 & yes \\
RA32source & 2024-25 & 356 & -2.77 & -2.792 & -0.02217 & yes \\
RA32source & 2025-26 & 528 & -2.644 & -2.619 & 0.02453 & yes \\
\bottomrule
\end{tabularx}
\par\smallskip
\begin{minipage}{\linewidth}\footnotesize\raggedright
Summary rows show valid/total folds and the number favoring affine in the last column; season rows show test n and convergence. Score is mean L/mu + log(mu), lower is better. Difference = power minus affine, so positive favors affine. Fold fits refit the training center. Training sets overlap; no iid-fold SE or four-season significance test is reported.
\end{minipage}
\end{table}

\begin{table}[htbp]
\centering
\scriptsize
\caption{Shared-count binomial thinning of the R\_\allowbreak{}A sample (n=1748): estimates, simultaneous bands and simulation diagnostics}\label{tab:thinning}
\setlength{\tabcolsep}{2pt}
\begin{tabularx}{\linewidth}{@{}>{\raggedright\arraybackslash}p{0.2400\linewidth}>{\raggedright\arraybackslash}p{0.1400\linewidth}>{\raggedright\arraybackslash}p{0.1400\linewidth}>{\raggedright\arraybackslash}p{0.1400\linewidth}>{\raggedright\arraybackslash}p{0.1400\linewidth}>{\raggedright\arraybackslash}p{0.1400\linewidth}@{}}
\toprule
Sample & q / fit & Metric & Estimate [95\% CI] & Simultaneous 95\% & Finite/B \\
\midrule
full & q=0.8 & b & 0.6747 [0.4905, 0.9034] & -- & 1000/1000 \\
full & q=0.8 & c & 0.4588 [0.3731, 0.5557] & [0.3381, 0.5795] & 1000/1000 \\
full & q=0.8 & delta\_\allowbreak{}a & -0.0003183 [-0.000786, 0.00005128] & [-0.0008833, 0.0002467] & 1000/1000 \\
full & q=0.8 & c\_\allowbreak{}over\_\allowbreak{}bP & 1.03 [0.8476, 1.242] & -- & 1000/1000 \\
full & q=0.6 & b & 0.8763 [0.6941, 1.103] & -- & 1000/1000 \\
full & q=0.6 & c & 0.4745 [0.4064, 0.5532] & [0.3751, 0.574] & 1000/1000 \\
full & q=0.6 & delta\_\allowbreak{}a & -0.000836 [-0.001963, -0.000006555] & [-0.00213, 0.0004583] & 1000/1000 \\
full & q=0.6 & c\_\allowbreak{}over\_\allowbreak{}bP & 1.065 [0.915, 1.241] & -- & 1000/1000 \\
full & q=0.5 & b & 1.057 [0.858, 1.275] & -- & 1000/1000 \\
full & q=0.5 & c & 0.4969 [0.4243, 0.5606] & [0.4047, 0.589] & 1000/1000 \\
full & q=0.5 & delta\_\allowbreak{}a & -0.001318 [-0.002744, -0.00009326] & [-0.003126, 0.0004903] & 1000/1000 \\
full & q=0.5 & c\_\allowbreak{}over\_\allowbreak{}bP & 1.115 [0.9621, 1.252] & -- & 1000/1000 \\
full & q=0.4 & b & 1.31 [1.112, 1.534] & -- & 1000/1000 \\
full & q=0.4 & c & 0.4997 [0.4397, 0.5611] & [0.4158, 0.5837] & 1000/1000 \\
full & q=0.4 & delta\_\allowbreak{}a & -0.001794 [-0.00376, -0.0001502] & [-0.004272, 0.0006837] & 1000/1000 \\
full & q=0.4 & c\_\allowbreak{}over\_\allowbreak{}bP & 1.122 [0.9903, 1.257] & -- & 1000/1000 \\
full & q=0.3 & b & 1.742 [1.503, 1.958] & -- & 1000/1000 \\
full & q=0.3 & c & 0.5067 [0.4482, 0.5539] & [0.4326, 0.5808] & 1000/1000 \\
full & q=0.3 & delta\_\allowbreak{}a & -0.002454 [-0.004852, -0.000159] & [-0.005728, 0.0008198] & 1000/1000 \\
full & q=0.3 & c\_\allowbreak{}over\_\allowbreak{}bP & 1.137 [1.008, 1.249] & -- & 1000/1000 \\
full & q=0.2 & b & 2.46 [2.221, 2.711] & -- & 1000/1000 \\
full & q=0.2 & c & 0.4749 [0.4292, 0.5229] & [0.4126, 0.5373] & 1000/1000 \\
full & q=0.2 & delta\_\allowbreak{}a & -0.002484 [-0.005751, 0.0005387] & [-0.006867, 0.001899] & 1000/1000 \\
full & q=0.2 & c\_\allowbreak{}over\_\allowbreak{}bP & 1.066 [0.9728, 1.173] & -- & 1000/1000 \\
full & q=0.1 & b & 3.983 [3.509, 4.379] & -- & 1000/1000 \\
full & q=0.1 & c & 0.3804 [0.3337, 0.4204] & [0.3234, 0.4373] & 1000/1000 \\
full & q=0.1 & delta\_\allowbreak{}a & 0.0007312 [-0.00305, 0.005984] & [-0.005595, 0.007057] & 1000/1000 \\
full & q=0.1 & c\_\allowbreak{}over\_\allowbreak{}bP & 0.8537 [0.7415, 0.9566] & -- & 1000/1000 \\
full & full & c & 0.4051 [0.3656, 0.4413] & -- & 1000/1000 \\
full & full & c\_\allowbreak{}over\_\allowbreak{}bP & 0.9093 [0.8101, 0.9995] & -- & 1000/1000 \\
full & full & delta\_\allowbreak{}a\_\allowbreak{}rate & -0.0001293 [-0.0006192, 0.0004699] & -- & 1000/1000 \\
full & noq01 & c & 0.4848 [0.4387, 0.5279] & -- & 1000/1000 \\
full & noq01 & c\_\allowbreak{}over\_\allowbreak{}bP & 1.088 [0.9936, 1.186] & -- & 1000/1000 \\
full & noq01 & delta\_\allowbreak{}a\_\allowbreak{}rate & -0.0008058 [-0.001703, 0.00005179] & -- & 1000/1000 \\
highCt50 & q=0.8 & b & 0.7797 [-0.007567, 1.573] & -- & 1000/1000 \\
highCt50 & q=0.8 & c & 0.442 [0.2771, 0.6691] & [0.1838, 0.7002] & 1000/1000 \\
highCt50 & q=0.8 & delta\_\allowbreak{}a & -0.00007119 [-0.0004223, 0.0001553] & [-0.0004703, 0.0003279] & 1000/1000 \\
highCt50 & q=0.8 & c\_\allowbreak{}over\_\allowbreak{}bP & 1.051 [0.6634, 1.606] & -- & 1000/1000 \\
highCt50 & q=0.6 & b & 0.969 [0.1971, 1.773] & -- & 1000/1000 \\
highCt50 & q=0.6 & c & 0.4497 [0.3339, 0.6318] & [0.256, 0.6433] & 1000/1000 \\
highCt50 & q=0.6 & delta\_\allowbreak{}a & -0.0001877 [-0.0009664, 0.0002652] & [-0.0009926, 0.0006171] & 1000/1000 \\
highCt50 & q=0.6 & c\_\allowbreak{}over\_\allowbreak{}bP & 1.069 [0.8015, 1.51] & -- & 1000/1000 \\
highCt50 & q=0.5 & b & 1.126 [0.3593, 1.891] & -- & 1000/1000 \\
highCt50 & q=0.5 & c & 0.4567 [0.3405, 0.6213] & [0.273, 0.6404] & 1000/1000 \\
highCt50 & q=0.5 & delta\_\allowbreak{}a & -0.0002913 [-0.001459, 0.0003529] & [-0.001456, 0.000873] & 1000/1000 \\
highCt50 & q=0.5 & c\_\allowbreak{}over\_\allowbreak{}bP & 1.086 [0.8187, 1.476] & -- & 1000/1000 \\
highCt50 & q=0.4 & b & 1.351 [0.6322, 2.134] & -- & 1000/1000 \\
highCt50 & q=0.4 & c & 0.4548 [0.3568, 0.593] & [0.2972, 0.6124] & 1000/1000 \\
highCt50 & q=0.4 & delta\_\allowbreak{}a & -0.0004078 [-0.001854, 0.0004131] & [-0.001938, 0.001122] & 1000/1000 \\
highCt50 & q=0.4 & c\_\allowbreak{}over\_\allowbreak{}bP & 1.081 [0.8483, 1.413] & -- & 1000/1000 \\
highCt50 & q=0.3 & b & 1.723 [1.005, 2.495] & -- & 1000/1000 \\
highCt50 & q=0.3 & c & 0.4515 [0.3668, 0.5711] & [0.3176, 0.5854] & 1000/1000 \\
highCt50 & q=0.3 & delta\_\allowbreak{}a & -0.0005649 [-0.00247, 0.0006027] & [-0.002578, 0.001448] & 1000/1000 \\
highCt50 & q=0.3 & c\_\allowbreak{}over\_\allowbreak{}bP & 1.073 [0.8687, 1.355] & -- & 1000/1000 \\
highCt50 & q=0.2 & b & 2.522 [1.77, 3.276] & -- & 1000/1000 \\
highCt50 & q=0.2 & c & 0.4633 [0.3813, 0.5611] & [0.3435, 0.5831] & 1000/1000 \\
highCt50 & q=0.2 & delta\_\allowbreak{}a & -0.0009427 [-0.00351, 0.0008025] & [-0.003822, 0.001937] & 1000/1000 \\
highCt50 & q=0.2 & c\_\allowbreak{}over\_\allowbreak{}bP & 1.102 [0.9074, 1.325] & -- & 1000/1000 \\
highCt50 & q=0.1 & b & 4.97 [4.216, 5.814] & -- & 999/1000 \\
highCt50 & q=0.1 & c & 0.4778 [0.4127, 0.5552] & [0.3835, 0.5721] & 999/1000 \\
highCt50 & q=0.1 & delta\_\allowbreak{}a & -0.001594 [-0.005615, 0.001213] & [-0.006149, 0.002962] & 999/1000 \\
highCt50 & q=0.1 & c\_\allowbreak{}over\_\allowbreak{}bP & 1.136 [0.9843, 1.323] & -- & 999/1000 \\
highCt50 & full & c & 0.4734 [0.409, 0.5503] & -- & 999/1000 \\
highCt50 & full & c\_\allowbreak{}over\_\allowbreak{}bP & 1.126 [0.9702, 1.316] & -- & 999/1000 \\
highCt50 & full & delta\_\allowbreak{}a\_\allowbreak{}rate & -0.0001928 [-0.0006891, 0.0001441] & -- & 999/1000 \\
highCt50 & noq01 & c & 0.4594 [0.3808, 0.5603] & -- & 1000/1000 \\
highCt50 & noq01 & c\_\allowbreak{}over\_\allowbreak{}bP & 1.092 [0.9026, 1.327] & -- & 1000/1000 \\
highCt50 & noq01 & delta\_\allowbreak{}a\_\allowbreak{}rate & -0.0002434 [-0.0009563, 0.0002083] & -- & 1000/1000 \\
both50 & q=0.8 & b & 0.5473 [-0.2021, 1.391] & -- & 1000/1000 \\
both50 & q=0.8 & c & 0.4994 [0.3073, 0.7827] & [0.1979, 0.8009] & 1000/1000 \\
both50 & q=0.8 & delta\_\allowbreak{}a & -0.0001674 [-0.0006684, 0.0001096] & [-0.0006864, 0.0003517] & 1000/1000 \\
both50 & q=0.8 & c\_\allowbreak{}over\_\allowbreak{}bP & 1.194 [0.722, 1.896] & -- & 1000/1000 \\
both50 & q=0.6 & b & 0.754 [0.02843, 1.586] & -- & 1000/1000 \\
both50 & q=0.6 & c & 0.4973 [0.3582, 0.7249] & [0.2629, 0.7317] & 1000/1000 \\
both50 & q=0.6 & delta\_\allowbreak{}a & -0.0004007 [-0.001555, 0.000186] & [-0.001503, 0.0007021] & 1000/1000 \\
both50 & q=0.6 & c\_\allowbreak{}over\_\allowbreak{}bP & 1.189 [0.854, 1.751] & -- & 1000/1000 \\
both50 & q=0.5 & b & 0.9288 [0.1947, 1.748] & -- & 1000/1000 \\
both50 & q=0.5 & c & 0.5063 [0.3635, 0.7058] & [0.2897, 0.723] & 1000/1000 \\
both50 & q=0.5 & delta\_\allowbreak{}a & -0.0006083 [-0.00219, 0.0001859] & [-0.002149, 0.0009326] & 1000/1000 \\
both50 & q=0.5 & c\_\allowbreak{}over\_\allowbreak{}bP & 1.21 [0.8863, 1.69] & -- & 1000/1000 \\
both50 & q=0.4 & b & 1.161 [0.4133, 1.977] & -- & 1000/1000 \\
both50 & q=0.4 & c & 0.4921 [0.3785, 0.6457] & [0.3095, 0.6746] & 1000/1000 \\
both50 & q=0.4 & delta\_\allowbreak{}a & -0.00078 [-0.002694, 0.0002525] & [-0.002738, 0.001178] & 1000/1000 \\
both50 & q=0.4 & c\_\allowbreak{}over\_\allowbreak{}bP & 1.176 [0.9064, 1.548] & -- & 1000/1000 \\
both50 & q=0.3 & b & 1.554 [0.8611, 2.363] & -- & 1000/1000 \\
both50 & q=0.3 & c & 0.4849 [0.3882, 0.6214] & [0.3325, 0.6373] & 1000/1000 \\
both50 & q=0.3 & delta\_\allowbreak{}a & -0.001058 [-0.003465, 0.0003471] & [-0.003573, 0.001457] & 1000/1000 \\
both50 & q=0.3 & c\_\allowbreak{}over\_\allowbreak{}bP & 1.159 [0.9271, 1.485] & -- & 1000/1000 \\
both50 & q=0.2 & b & 2.382 [1.63, 3.167] & -- & 1000/1000 \\
both50 & q=0.2 & c & 0.4899 [0.3955, 0.5995] & [0.3589, 0.6209] & 1000/1000 \\
both50 & q=0.2 & delta\_\allowbreak{}a & -0.001588 [-0.004831, 0.0005575] & [-0.005049, 0.001874] & 1000/1000 \\
both50 & q=0.2 & c\_\allowbreak{}over\_\allowbreak{}bP & 1.171 [0.9456, 1.432] & -- & 1000/1000 \\
both50 & q=0.1 & b & 4.863 [4.1, 5.734] & -- & 999/1000 \\
both50 & q=0.1 & c & 0.4934 [0.4221, 0.5729] & [0.3956, 0.5913] & 999/1000 \\
both50 & q=0.1 & delta\_\allowbreak{}a & -0.002397 [-0.007015, 0.000824] & [-0.007548, 0.002755] & 999/1000 \\
both50 & q=0.1 & c\_\allowbreak{}over\_\allowbreak{}bP & 1.18 [1.005, 1.371] & -- & 999/1000 \\
both50 & full & c & 0.4926 [0.42, 0.5742] & -- & 999/1000 \\
both50 & full & c\_\allowbreak{}over\_\allowbreak{}bP & 1.178 [1.004, 1.379] & -- & 999/1000 \\
both50 & full & delta\_\allowbreak{}a\_\allowbreak{}rate & -0.0003058 [-0.000896, 0.00009878] & -- & 999/1000 \\
both50 & noq01 & c & 0.4898 [0.3971, 0.6072] & -- & 1000/1000 \\
both50 & noq01 & c\_\allowbreak{}over\_\allowbreak{}bP & 1.171 [0.9533, 1.462] & -- & 1000/1000 \\
both50 & noq01 & delta\_\allowbreak{}a\_\allowbreak{}rate & -0.0004328 [-0.001331, 0.0001204] & -- & 1000/1000 \\
\bottomrule
\end{tabularx}
\par\smallskip
\begin{minipage}{\linewidth}\footnotesize\raggedright
Per-q intervals are pointwise; the separate bands are simultaneous over q within each subset. full uses all q; noq01 excludes q=0.1. highCt50 requires original current count >=50; both50 requires both original adjacent counts >=50. c is the slope increment per (1/q-1); bP is the plug-in Poisson reference. B=1000; R(point)=400; R(inner, actual)=50; precision criterion: MC SE at most 0.1 times the point estimate, applied to the five monitored response summaries; max finite MC SE/|point| over all stored metrics=68.44; max finite MC SE/sampling SE over all stored metrics=1.74; undefined MC ratios=0/219; MC status=tolerance\_\allowbreak{}met. Point fits 8400/8400; bootstrap fits 1049998/1050000; failed bootstrap averages=2. Any failed inner fit makes that q average undefined; no silent averaging over successful fits. full: n=1748; bP=0.4455 [0.4282, 0.4658]. highCt50: n=828; bP=0.4206 [0.4032, 0.4387]. both50: n=821; bP=0.4183 [0.4008, 0.4359].
\end{minipage}
\end{table}

\begin{table}[htbp]
\centering
\scriptsize
\caption{Fixed effects, paired ratios, block sensitivity and convergence}\label{tab:fe}
\setlength{\tabcolsep}{2pt}
\begin{tabularx}{\linewidth}{@{}>{\raggedright\arraybackslash}p{0.2400\linewidth}>{\raggedright\arraybackslash}p{0.1750\linewidth}>{\raggedright\arraybackslash}p{0.1750\linewidth}>{\raggedright\arraybackslash}p{0.1750\linewidth}>{\raggedright\arraybackslash}p{0.1750\linewidth}@{}}
\toprule
Block & Parameter & Estimate [95\% CI] & Finite/B & Nonpositive / diagnostic \\
\midrule
L=1 & Pooled a & 0.01495 [0.009976, 0.02114] & 1000/1000 & 0 \\
L=1 & Pooled b & 0.56 [0.3682, 0.7707] & 1000/1000 & 0 \\
L=1 & Pooled b/a & 37.44 [18.64, 73.92] & 1000/1000 & 0 \\
L=1 & State FE b & 0.3794 [0.1082, 0.6403] & 999/1000 & 0 \\
L=1 & State FE mean a & 0.02192 [0.01289, 0.03502] & 999/1000 & 0 \\
L=1 & State FE b/mean a & 17.3 [3.315, 47.8] & 999/1000 & 0 \\
L=1 & State FE median a & 0.01751 [0.01037, 0.03014] & 999/1000 & 0 \\
L=1 & State FE b/median a & 21.67 [3.977, 57.85] & 999/1000 & 0 \\
L=1 & State FE row-weighted a & 0.02061 [0.01339, 0.03114] & 999/1000 & 0 \\
L=1 & State FE row-weighted ratio & 18.41 [3.705, 45.56] & 999/1000 & 0 \\
L=1 & State FE count-weighted a & 0.01542 [0.01056, 0.02288] & 999/1000 & 0 \\
L=1 & State FE count-weighted ratio & 24.61 [5.265, 55.65] & 999/1000 & 0 \\
L=1 & True D\_\allowbreak{}t b & 0.3691 [0.1153, 0.6145] & 999/1000 & 0 \\
L=1 & True D\_\allowbreak{}t mean a & 0.02084 [0.01218, 0.03429] & 999/1000 & 0 \\
L=1 & True D\_\allowbreak{}t ratio & 17.71 [3.267, 47.81] & 999/1000 & 0 \\
L=1 & Paired pooled/FE ratio & 2.164 [1.04, 7.745] & 999/1000 & 0 \\
L=2 & Pooled a & 0.01495 [0.009493, 0.02077] & 1000/1000 & 0 \\
L=2 & Pooled b & 0.56 [0.371, 0.763] & 1000/1000 & 0 \\
L=2 & Pooled b/a & 37.44 [18.73, 74.77] & 1000/1000 & 0 \\
L=2 & State FE b & 0.3794 [0.109, 0.6611] & 1000/1000 & 0 \\
L=2 & State FE mean a & 0.02192 [0.01288, 0.03581] & 1000/1000 & 0 \\
L=2 & State FE b/mean a & 17.3 [3.453, 50.16] & 1000/1000 & 0 \\
L=2 & State FE median a & 0.01751 [0.01024, 0.03013] & 1000/1000 & 0 \\
L=2 & State FE b/median a & 21.67 [3.984, 60.15] & 1000/1000 & 0 \\
L=2 & State FE row-weighted a & 0.02061 [0.013, 0.03107] & 1000/1000 & 0 \\
L=2 & State FE row-weighted ratio & 18.41 [3.819, 48.17] & 1000/1000 & 0 \\
L=2 & State FE count-weighted a & 0.01542 [0.0102, 0.02293] & 1000/1000 & 0 \\
L=2 & State FE count-weighted ratio & 24.61 [5.111, 57.08] & 1000/1000 & 0 \\
L=2 & True D\_\allowbreak{}t b & 0.3691 [0.1082, 0.612] & 1000/1000 & 0.001 \\
L=2 & True D\_\allowbreak{}t mean a & 0.02084 [0.01168, 0.03555] & 1000/1000 & 0 \\
L=2 & True D\_\allowbreak{}t ratio & 17.71 [3.26, 51.13] & 1000/1000 & 0.001 \\
L=2 & Paired pooled/FE ratio & 2.164 [0.964, 8.483] & 1000/1000 & 0 \\
L=3 & Pooled a & 0.01495 [0.01001, 0.02089] & 1000/1000 & 0 \\
L=3 & Pooled b & 0.56 [0.3851, 0.7603] & 1000/1000 & 0 \\
L=3 & Pooled b/a & 37.44 [19.39, 68.59] & 1000/1000 & 0 \\
L=3 & State FE b & 0.3794 [0.111, 0.6669] & 1000/1000 & 0.001 \\
L=3 & State FE mean a & 0.02192 [0.01247, 0.03692] & 1000/1000 & 0 \\
L=3 & State FE b/mean a & 17.3 [3.051, 50.11] & 1000/1000 & 0.001 \\
L=3 & State FE median a & 0.01751 [0.01068, 0.03045] & 1000/1000 & 0 \\
L=3 & State FE b/median a & 21.67 [3.982, 58.24] & 1000/1000 & 0.001 \\
L=3 & State FE row-weighted a & 0.02061 [0.01293, 0.0317] & 1000/1000 & 0 \\
L=3 & State FE row-weighted ratio & 18.41 [3.407, 47.65] & 1000/1000 & 0.001 \\
L=3 & State FE count-weighted a & 0.01542 [0.01076, 0.02338] & 1000/1000 & 0 \\
L=3 & State FE count-weighted ratio & 24.61 [4.841, 57.62] & 1000/1000 & 0.001 \\
L=3 & True D\_\allowbreak{}t b & 0.3691 [0.1072, 0.6252] & 1000/1000 & 0.001 \\
L=3 & True D\_\allowbreak{}t mean a & 0.02084 [0.0116, 0.03629] & 1000/1000 & 0 \\
L=3 & True D\_\allowbreak{}t ratio & 17.71 [3.068, 51.41] & 1000/1000 & 0.001 \\
L=3 & Paired pooled/FE ratio & 2.164 [0.9866, 8.87] & 999/1000 & 0 \\
Point diagnostic & week\_\allowbreak{}only & b=0.5258 & No CI & group\_\allowbreak{}gamma\_\allowbreak{}objective\_\allowbreak{}fallback \\
Point diagnostic & state\_\allowbreak{}and\_\allowbreak{}week & Unavailable & No CI & iteration\_\allowbreak{}limit \\
\bottomrule
\end{tabularx}
\par\smallskip
\begin{minipage}{\linewidth}\footnotesize\raggedright
Intervals retain signed draws and are conditional on finite estimates when finite/B < 1. Last column gives the nonpositive fraction over all requested draws. FE targets the sampled state occurrences; pooled/FE comparisons are paired within draws. L denotes the within-season week-block length. Week-only and state-and-week fits are point diagnostics; failed fits have no estimate or interval. L=1 pooled: 1000 converged, 0 failed; \{\}. L=1 state: 999 converged, 1 failed; \{"gamma\_\allowbreak{}objective\_\allowbreak{}fallback\_\allowbreak{}failed": 1\}. L=1 state\_\allowbreak{}D: 999 converged, 1 failed; \{"gamma\_\allowbreak{}objective\_\allowbreak{}fallback\_\allowbreak{}failed": 1\}. L=2 pooled: 1000 converged, 0 failed; \{\}. L=2 state: 1000 converged, 0 failed; \{\}. L=2 state\_\allowbreak{}D: 1000 converged, 0 failed; \{\}. L=3 pooled: 1000 converged, 0 failed; \{\}. L=3 state: 1000 converged, 0 failed; \{\}. L=3 state\_\allowbreak{}D: 1000 converged, 0 failed; \{\}.
\end{minipage}
\end{table}

\begin{table}[htbp]
\centering
\scriptsize
\caption{Centering stages with a common loss normalizer}\label{tab:centering}
\setlength{\tabcolsep}{2pt}
\begin{tabularx}{\linewidth}{@{}>{\raggedright\arraybackslash}p{0.2400\linewidth}>{\raggedright\arraybackslash}p{0.1750\linewidth}>{\raggedright\arraybackslash}p{0.1750\linewidth}>{\raggedright\arraybackslash}p{0.1750\linewidth}>{\raggedright\arraybackslash}p{0.1750\linewidth}@{}}
\toprule
Statistic / center & Pooled mean & a / change [95\% CI] & b [95\% CI] & Finite/B \\
\midrule
RA / pooled & 1.206 & 0.01495 [0.01016, 0.02083] & 0.56 [0.3625, 0.7639] & 1000/1000 \\
RA / season & 1.206 & 0.01471 [0.009507, 0.01966] & 0.5391 [0.3602, 0.7205] & 1000/1000 \\
RA / season\_\allowbreak{}linear & 1.206 & 0.0154 [0.008972, 0.01901] & 0.3936 [0.2717, 0.5763] & 1000/1000 \\
RA / week & 1.206 & 0.009766 [0.007103, 0.01357] & 0.4581 [0.3124, 0.5916] & 1000/1000 \\
RA / week\_\allowbreak{}leave\_\allowbreak{}state\_\allowbreak{}out & 1.206 & 0.01016 [0.007654, 0.01481] & 0.4841 [0.3492, 0.6636] & 1000/1000 \\
RA / week\_\allowbreak{}change & 1.206 & -0.005189 [-0.01015, -0.0006938] & -- & 1000/1000 \\
RA / week\_\allowbreak{}reduction\_\allowbreak{}fraction & 1.206 & 34.7 [5.967, 51.9] & -- & 1000/1000 \\
G / pooled & 1.724 & 0.08401 [0.05163, 0.1222] & 1.971 [0.9987, 2.986] & 1000/1000 \\
G / season & 1.724 & 0.08163 [0.0504, 0.1156] & 1.842 [0.9591, 2.602] & 1000/1000 \\
G / season\_\allowbreak{}linear & 1.724 & 0.07864 [0.04808, 0.09897] & 1.249 [0.6423, 1.988] & 1000/1000 \\
G / week & 1.724 & 0.05607 [0.03962, 0.07812] & 1.401 [0.7134, 1.912] & 1000/1000 \\
G / week\_\allowbreak{}leave\_\allowbreak{}state\_\allowbreak{}out & 1.724 & 0.05843 [0.04308, 0.08505] & 1.477 [0.8089, 2.138] & 1000/1000 \\
G / week\_\allowbreak{}change & 1.724 & -0.02794 [-0.05687, -0.003708] & -- & 1000/1000 \\
G / week\_\allowbreak{}reduction\_\allowbreak{}fraction & 1.724 & 33.25 [6.805, 51.63] & -- & 1000/1000 \\
\bottomrule
\end{tabularx}
\par\smallskip
\begin{minipage}{\linewidth}\footnotesize\raggedright
All five centers use the same pooled mean as the relative-loss denominator within each statistic; the denominator is refitted with weights within each bootstrap draw. week\_\allowbreak{}leave\_\allowbreak{}state\_\allowbreak{}out removes the focal state from its weekly center. Reduction is 100*(a\_\allowbreak{}pooled-a\_\allowbreak{}week)/a\_\allowbreak{}pooled, in percent; changes are descriptive and not additive variance components. RA/pooled: 1000/1000 finite draws; RA/season: 1000/1000 finite draws; RA/season\_\allowbreak{}linear: 1000/1000 finite draws; RA/week: 1000/1000 finite draws; RA/week\_\allowbreak{}leave\_\allowbreak{}state\_\allowbreak{}out: 1000/1000 finite draws; G/pooled: 1000/1000 finite draws; G/season: 1000/1000 finite draws; G/season\_\allowbreak{}linear: 1000/1000 finite draws; G/week: 1000/1000 finite draws; G/week\_\allowbreak{}leave\_\allowbreak{}state\_\allowbreak{}out: 1000/1000 finite draws.
\end{minipage}
\end{table}

\begin{table}[htbp]
\centering
\scriptsize
\caption{Conditional calibration across all declared scenarios}\label{tab:calibration}
\setlength{\tabcolsep}{2pt}
\begin{tabularx}{\linewidth}{@{}>{\raggedright\arraybackslash}p{0.2400\linewidth}>{\raggedright\arraybackslash}p{0.1400\linewidth}>{\raggedright\arraybackslash}p{0.1400\linewidth}>{\raggedright\arraybackslash}p{0.1400\linewidth}>{\raggedright\arraybackslash}p{0.1400\linewidth}>{\raggedright\arraybackslash}p{0.1400\linewidth}@{}}
\toprule
Scenario / statistic & v [95\% CI] & Finite/B & Nonpositive / nonfinite & Outside / slope <=0 & Failed fits / all \\
\midrule
delay\_\allowbreak{}long / RA & 0.04587 [0.02996, 0.06461] & 2000/2000 & 0/0 & 0/0 & 0/1000 \\
delay\_\allowbreak{}long / G & 0.04686 [0.02666, 0.06952] & 2000/2000 & 0/0 & 3/0 & 0/1000 \\
delay\_\allowbreak{}short / RA & 0.03863 [0.02531, 0.05527] & 2000/2000 & 0/0 & 0/0 & 0/1000 \\
delay\_\allowbreak{}short / G & 0.04161 [0.02382, 0.06314] & 2000/2000 & 0/0 & 0/0 & 0/1000 \\
joint\_\allowbreak{}gi3.2\_\allowbreak{}pn0.25\_\allowbreak{}obs1 / RA & 0.02882 [0.01937, 0.04107] & 2000/2000 & 0/0 & 0/0 & 0/1000 \\
joint\_\allowbreak{}gi3.2\_\allowbreak{}pn0.25\_\allowbreak{}obs1 / G & 0.02967 [0.01709, 0.04456] & 2000/2000 & 0/0 & 0/0 & 0/1000 \\
joint\_\allowbreak{}gi3.2\_\allowbreak{}pn0.25\_\allowbreak{}obs1.25 / RA & 0.02844 [0.01853, 0.03995] & 2000/2000 & 0/0 & 0/0 & 0/1000 \\
joint\_\allowbreak{}gi3.2\_\allowbreak{}pn0.25\_\allowbreak{}obs1.25 / G & 0.0294 [0.01681, 0.04331] & 2000/2000 & 0/0 & 0/0 & 0/1000 \\
joint\_\allowbreak{}gi3.2\_\allowbreak{}pn0\_\allowbreak{}obs1 / RA & 0.03178 [0.02085, 0.04488] & 2000/2000 & 0/0 & 0/0 & 0/1000 \\
joint\_\allowbreak{}gi3.2\_\allowbreak{}pn0\_\allowbreak{}obs1 / G & 0.03369 [0.01924, 0.05045] & 2000/2000 & 0/0 & 0/0 & 0/1000 \\
joint\_\allowbreak{}gi3.2\_\allowbreak{}pn0\_\allowbreak{}obs1.25 / RA & 0.03128 [0.02019, 0.04538] & 2000/2000 & 0/0 & 0/0 & 0/1000 \\
joint\_\allowbreak{}gi3.2\_\allowbreak{}pn0\_\allowbreak{}obs1.25 / G & 0.03348 [0.01895, 0.05184] & 2000/2000 & 0/0 & 0/0 & 0/1000 \\
joint\_\allowbreak{}gi3.6\_\allowbreak{}pn0.25\_\allowbreak{}obs1 / RA & 0.03764 [0.02444, 0.05359] & 2000/2000 & 0/0 & 0/0 & 0/1000 \\
joint\_\allowbreak{}gi3.6\_\allowbreak{}pn0.25\_\allowbreak{}obs1 / G & 0.03902 [0.02247, 0.05689] & 2000/2000 & 0/0 & 0/0 & 0/1000 \\
joint\_\allowbreak{}gi3.6\_\allowbreak{}pn0.25\_\allowbreak{}obs1.25 / RA & 0.03638 [0.02354, 0.0519] & 2000/2000 & 0/0 & 0/0 & 0/1000 \\
joint\_\allowbreak{}gi3.6\_\allowbreak{}pn0.25\_\allowbreak{}obs1.25 / G & 0.03785 [0.0216, 0.05704] & 2000/2000 & 0/0 & 0/0 & 0/1000 \\
joint\_\allowbreak{}gi3.6\_\allowbreak{}pn0\_\allowbreak{}obs1 / RA & 0.04002 [0.02558, 0.05767] & 2000/2000 & 0/0 & 0/0 & 0/1000 \\
joint\_\allowbreak{}gi3.6\_\allowbreak{}pn0\_\allowbreak{}obs1 / G & 0.04282 [0.02407, 0.06368] & 2000/2000 & 0/0 & 0/0 & 0/1000 \\
joint\_\allowbreak{}gi3.6\_\allowbreak{}pn0\_\allowbreak{}obs1.25 / RA & 0.04035 [0.02648, 0.0569] & 2000/2000 & 0/0 & 0/0 & 0/1000 \\
joint\_\allowbreak{}gi3.6\_\allowbreak{}pn0\_\allowbreak{}obs1.25 / G & 0.04356 [0.02576, 0.064] & 2000/2000 & 0/0 & 0/0 & 0/1000 \\
k01 / RA & 0.04023 [0.02553, 0.05736] & 2000/2000 & 0/0 & 0/0 & 0/1000 \\
k01 / G & 0.04289 [0.02475, 0.0645] & 2000/2000 & 0/0 & 0/0 & 0/1000 \\
matched32\_\allowbreak{}binomial / RA & 0.03108 [0.02037, 0.04399] & 2000/2000 & 0/0 & 0/0 & 0/1000 \\
matched32\_\allowbreak{}binomial / G & 0.03287 [0.0186, 0.0489] & 2000/2000 & 0/0 & 0/0 & 0/1000 \\
matched32\_\allowbreak{}source\_\allowbreak{}sd21\_\allowbreak{}binomial / RA & 0.02884 [0.01908, 0.04027] & 2000/2000 & 0/0 & 0/0 & 0/1000 \\
matched32\_\allowbreak{}source\_\allowbreak{}sd21\_\allowbreak{}binomial / G & 0.031 [0.01794, 0.04533] & 2000/2000 & 0/0 & 0/0 & 0/1000 \\
matched36\_\allowbreak{}binomial / RA & 0.04 [0.02623, 0.0562] & 2000/2000 & 0/0 & 0/0 & 0/1000 \\
matched36\_\allowbreak{}binomial / G & 0.04298 [0.02466, 0.06284] & 2000/2000 & 0/0 & 0/0 & 0/1000 \\
national\_\allowbreak{}persistence25 / RA & 0.0369 [0.0238, 0.05109] & 2000/2000 & 0/0 & 0/0 & 0/1000 \\
national\_\allowbreak{}persistence25 / G & 0.03817 [0.02185, 0.05643] & 2000/2000 & 0/0 & 0/0 & 0/1000 \\
national\_\allowbreak{}share0 / RA & 0.03935 [0.0255, 0.05576] & 2000/2000 & 0/0 & 0/0 & 0/1000 \\
national\_\allowbreak{}share0 / G & 0.0416 [0.02393, 0.06162] & 2000/2000 & 0/0 & 0/0 & 0/1000 \\
national\_\allowbreak{}share65 / RA & 0.04014 [0.02576, 0.05685] & 2000/2000 & 0/0 & 0/0 & 0/1000 \\
national\_\allowbreak{}share65 / G & 0.04327 [0.02461, 0.06553] & 2000/2000 & 0/0 & 0/0 & 0/1000 \\
reporting\_\allowbreak{}shock / RA & 0.01954 [0.006213, 0.0358] & 2000/2000 & 0/0 & 0/0 & 0/1000 \\
reporting\_\allowbreak{}shock / G & 0.02903 [0.0117, 0.04918] & 2000/2000 & 0/0 & 0/0 & 0/1000 \\
rho\_\allowbreak{}h005 / RA & 0.03991 [0.02639, 0.05628] & 2000/2000 & 0/0 & 0/0 & 0/1000 \\
rho\_\allowbreak{}h005 / G & 0.04283 [0.02548, 0.06297] & 2000/2000 & 0/0 & 0/0 & 0/1000 \\
rho\_\allowbreak{}h02 / RA & 0.04029 [0.02632, 0.05694] & 2000/2000 & 0/0 & 0/0 & 0/1000 \\
rho\_\allowbreak{}h02 / G & 0.04317 [0.02494, 0.06466] & 2000/2000 & 0/0 & 0/0 & 0/1000 \\
state\_\allowbreak{}persistence50 / RA & 0.03244 [0.02095, 0.04559] & 2000/2000 & 0/0 & 0/0 & 0/1000 \\
state\_\allowbreak{}persistence50 / G & 0.03137 [0.01788, 0.0467] & 2000/2000 & 0/0 & 0/0 & 0/1000 \\
true2.2\_\allowbreak{}analyst36 / RA & 0.0146 [0.009442, 0.02045] & 2000/2000 & 0/0 & 0/0 & 0/1000 \\
true2.2\_\allowbreak{}analyst36 / G & 0.01473 [0.008358, 0.02141] & 2000/2000 & 0/0 & 0/0 & 0/1000 \\
true2.9\_\allowbreak{}analyst36 / RA & 0.02551 [0.0168, 0.03572] & 2000/2000 & 0/0 & 0/0 & 0/1000 \\
true2.9\_\allowbreak{}analyst36 / G & 0.02679 [0.01524, 0.0401] & 2000/2000 & 0/0 & 0/0 & 0/1000 \\
true4.3\_\allowbreak{}analyst36 / RA & 0.05807 [0.03763, 0.08382] & 2000/2000 & 0/0 & 27/0 & 0/1000 \\
true4.3\_\allowbreak{}analyst36 / G & 0.06312 [0.0366, 0.0955] & 2000/2000 & 0/0 & 78/0 & 0/1000 \\
\bottomrule
\end{tabularx}
\par\smallskip
\begin{minipage}{\linewidth}\footnotesize\raggedright
v is a model-conditional weekly marginal variance before the lower bound, not an identified biological variance or upper bound. Intervals are the 2.5th and 97.5th percentiles of all finite signed draws; no positive-only filtering or clamping is applied. Columns report nonpositive/nonfinite inversion draws, out-of-grid/nonpositive-slope draws, and failed/attempted simulation fits. delay\_\allowbreak{}long: grid=[0.0, 0.02, 0.04, 0.06, 0.09]; RA point outside=False, nonpositive=False, max |line residual|/MC SE=0.8913; G point outside=False, nonpositive=False, max |line residual|/MC SE=0.5474. delay\_\allowbreak{}short: grid=[0.0, 0.02, 0.04, 0.06, 0.09]; RA point outside=False, nonpositive=False, max |line residual|/MC SE=1.781; G point outside=False, nonpositive=False, max |line residual|/MC SE=1.162. joint\_\allowbreak{}gi3.2\_\allowbreak{}pn0.25\_\allowbreak{}obs1: grid=[0.0, 0.02, 0.04, 0.06, 0.09]; RA point outside=False, nonpositive=False, max |line residual|/MC SE=0.974; G point outside=False, nonpositive=False, max |line residual|/MC SE=1.154. joint\_\allowbreak{}gi3.2\_\allowbreak{}pn0.25\_\allowbreak{}obs1.25: grid=[0.0, 0.02, 0.04, 0.06, 0.09]; RA point outside=False, nonpositive=False, max |line residual|/MC SE=1.557; G point outside=False, nonpositive=False, max |line residual|/MC SE=1.668. joint\_\allowbreak{}gi3.2\_\allowbreak{}pn0\_\allowbreak{}obs1: grid=[0.0, 0.02, 0.04, 0.06, 0.09]; RA point outside=False, nonpositive=False, max |line residual|/MC SE=1.13; G point outside=False, nonpositive=False, max |line residual|/MC SE=1.116. joint\_\allowbreak{}gi3.2\_\allowbreak{}pn0\_\allowbreak{}obs1.25: grid=[0.0, 0.02, 0.04, 0.06, 0.09]; RA point outside=False, nonpositive=False, max |line residual|/MC SE=1.369; G point outside=False, nonpositive=False, max |line residual|/MC SE=1.917. joint\_\allowbreak{}gi3.6\_\allowbreak{}pn0.25\_\allowbreak{}obs1: grid=[0.0, 0.02, 0.04, 0.06, 0.09]; RA point outside=False, nonpositive=False, max |line residual|/MC SE=1.666; G point outside=False, nonpositive=False, max |line residual|/MC SE=1.308. joint\_\allowbreak{}gi3.6\_\allowbreak{}pn0.25\_\allowbreak{}obs1.25: grid=[0.0, 0.02, 0.04, 0.06, 0.09]; RA point outside=False, nonpositive=False, max |line residual|/MC SE=1.255; G point outside=False, nonpositive=False, max |line residual|/MC SE=1.081. joint\_\allowbreak{}gi3.6\_\allowbreak{}pn0\_\allowbreak{}obs1: grid=[0.0, 0.02, 0.04, 0.06, 0.09]; RA point outside=False, nonpositive=False, max |line residual|/MC SE=1.596; G point outside=False, nonpositive=False, max |line residual|/MC SE=1.363. joint\_\allowbreak{}gi3.6\_\allowbreak{}pn0\_\allowbreak{}obs1.25: grid=[0.0, 0.02, 0.04, 0.06, 0.09]; RA point outside=False, nonpositive=False, max |line residual|/MC SE=0.9047; G point outside=False, nonpositive=False, max |line residual|/MC SE=0.9728. k01: grid=[0.0, 0.02, 0.04, 0.06, 0.09]; RA point outside=False, nonpositive=False, max |line residual|/MC SE=1.074; G point outside=False, nonpositive=False, max |line residual|/MC SE=1.429. matched32\_\allowbreak{}binomial: grid=[0.0, 0.02, 0.04, 0.06, 0.09]; RA point outside=False, nonpositive=False, max |line residual|/MC SE=0.4468; G point outside=False, nonpositive=False, max |line residual|/MC SE=0.4955. matched32\_\allowbreak{}source\_\allowbreak{}sd21\_\allowbreak{}binomial: grid=[0.0, 0.02, 0.04, 0.06, 0.09]; RA point outside=False, nonpositive=False, max |line residual|/MC SE=1.05; G point outside=False, nonpositive=False, max |line residual|/MC SE=0.9843. matched36\_\allowbreak{}binomial: grid=[0.0, 0.02, 0.04, 0.06, 0.09]; RA point outside=False, nonpositive=False, max |line residual|/MC SE=2.388; G point outside=False, nonpositive=False, max |line residual|/MC SE=1.795. national\_\allowbreak{}persistence25: grid=[0.0, 0.02, 0.04, 0.06, 0.09]; RA point outside=False, nonpositive=False, max |line residual|/MC SE=2.246; G point outside=False, nonpositive=False, max |line residual|/MC SE=2.37. national\_\allowbreak{}share0: grid=[0.0, 0.02, 0.04, 0.06, 0.09]; RA point outside=False, nonpositive=False, max |line residual|/MC SE=0.8455; G point outside=False, nonpositive=False, max |line residual|/MC SE=0.6427. national\_\allowbreak{}share65: grid=[0.0, 0.02, 0.04, 0.06, 0.09]; RA point outside=False, nonpositive=False, max |line residual|/MC SE=1.134; G point outside=False, nonpositive=False, max |line residual|/MC SE=0.7098. reporting\_\allowbreak{}shock: grid=[0.0, 0.02, 0.04, 0.06, 0.09]; RA point outside=False, nonpositive=False, max |line residual|/MC SE=0.8369; G point outside=False, nonpositive=False, max |line residual|/MC SE=0.9846. rho\_\allowbreak{}h005: grid=[0.0, 0.02, 0.04, 0.06, 0.09]; RA point outside=False, nonpositive=False, max |line residual|/MC SE=1.481; G point outside=False, nonpositive=False, max |line residual|/MC SE=1.628. rho\_\allowbreak{}h02: grid=[0.0, 0.02, 0.04, 0.06, 0.09]; RA point outside=False, nonpositive=False, max |line residual|/MC SE=1.184; G point outside=False, nonpositive=False, max |line residual|/MC SE=1.261. state\_\allowbreak{}persistence50: grid=[0.0, 0.02, 0.04, 0.06, 0.09]; RA point outside=False, nonpositive=False, max |line residual|/MC SE=3.259; G point outside=False, nonpositive=False, max |line residual|/MC SE=3.007. true2.2\_\allowbreak{}analyst36: grid=[0.0, 0.02, 0.04, 0.06, 0.09]; RA point outside=False, nonpositive=False, max |line residual|/MC SE=1.281; G point outside=False, nonpositive=False, max |line residual|/MC SE=1.023. true2.9\_\allowbreak{}analyst36: grid=[0.0, 0.02, 0.04, 0.06, 0.09]; RA point outside=False, nonpositive=False, max |line residual|/MC SE=0.5604; G point outside=False, nonpositive=False, max |line residual|/MC SE=1.366. true4.3\_\allowbreak{}analyst36: grid=[0.0, 0.02, 0.04, 0.06, 0.09]; RA point outside=False, nonpositive=False, max |line residual|/MC SE=1.162; G point outside=False, nonpositive=False, max |line residual|/MC SE=1.198.
\end{minipage}
\end{table}

\begin{table}[htbp]
\centering
\small
\caption{Model-conditional design size and power}\label{tab:design}
\setlength{\tabcolsep}{2pt}
\begin{tabularx}{\linewidth}{@{}>{\raggedright\arraybackslash}p{0.2400\linewidth}>{\raggedright\arraybackslash}p{0.1400\linewidth}>{\raggedright\arraybackslash}p{0.1400\linewidth}>{\raggedright\arraybackslash}p{0.1400\linewidth}>{\raggedright\arraybackslash}p{0.1400\linewidth}>{\raggedright\arraybackslash}p{0.1400\linewidth}@{}}
\toprule
Seasons & Count multiplier & Size [95\% CI] & Power [95\% CI] & Valid C/N/A & Size compatible \\
\midrule
4 & 1 & 0.058 [0.03918, 0.08224] & 0.996 [0.9856, 0.9995] & 500/500/500 & yes \\
4 & 2 & 0.06 [0.04084, 0.08455] & 0.998 [0.9889, 0.9999] & 500/500/500 & yes \\
4 & 4 & 0.052 [0.03425, 0.07527] & 0.994 [0.9826, 0.9988] & 500/500/500 & yes \\
8 & 1 & 0.058 [0.03918, 0.08224] & 1 [0.9926, 1] & 500/500/500 & yes \\
8 & 2 & 0.038 [0.02303, 0.05871] & 1 [0.9926, 1] & 500/500/500 & yes \\
8 & 4 & 0.082 [0.05949, 0.1096] & 1 [0.9926, 1] & 500/500/500 & no \\
12 & 1 & 0.058 [0.03918, 0.08224] & 1 [0.9926, 1] & 500/500/500 & yes \\
12 & 2 & 0.07 [0.04924, 0.09601] & 1 [0.9926, 1] & 500/500/500 & yes \\
12 & 4 & 0.054 [0.03588, 0.0776] & 1 [0.9926, 1] & 500/500/500 & yes \\
20 & 1 & 0.032 [0.0184, 0.05145] & 1 [0.9926, 1] & 500/500/500 & yes \\
20 & 2 & 0.058 [0.03918, 0.08224] & 1 [0.9926, 1] & 500/500/500 & yes \\
20 & 4 & 0.056 [0.03753, 0.07992] & 1 [0.9926, 1] & 500/500/500 & yes \\
\bottomrule
\end{tabularx}
\par\smallskip
\begin{minipage}{\linewidth}\footnotesize\raggedright
Conditional discrimination simulations recycle the observed season/count design. Calibration and evaluation samples are independent. Valid counts are calibration/null/alternative. The final column states whether the size interval contains nominal alpha; all cells are reported, including failed size checks. These simulations do not imply a universal season threshold or a universal 80\% power claim. alpha=0.05; replicates/stage=500.
\end{minipage}
\end{table}

\begin{table}[htbp]
\centering
\scriptsize
\caption{Season-specific fits and reporting participation sensitivity}\label{tab:secondary}
\setlength{\tabcolsep}{2pt}
\begin{tabularx}{\linewidth}{@{}>{\raggedright\arraybackslash}p{0.2400\linewidth}>{\raggedright\arraybackslash}p{0.1400\linewidth}>{\raggedright\arraybackslash}p{0.1400\linewidth}>{\raggedright\arraybackslash}p{0.1400\linewidth}>{\raggedright\arraybackslash}p{0.1400\linewidth}>{\raggedright\arraybackslash}p{0.1400\linewidth}@{}}
\toprule
Statistic / period & n & a / change [95\% CI] & b / original a & b/a / adjusted a & Finite/B \\
\midrule
RA / 2022-23 & 365 & 0.01369 [0.0071, 0.02102] & 0.6967 & 50.91 & 1000/1000 \\
RA / 2023-24 & 499 & 0.006174 [0.002667, 0.01542] & 0.7018 & 113.7 & 1000/1000 \\
RA / 2024-25 & 356 & 0.01268 [0.005687, 0.02276] & 0.5826 & 45.93 & 1000/1000 \\
RA / 2025-26 & 528 & 0.02253 [0.008432, 0.03509] & 0.3533 & 15.68 & 1000/1000 \\
RA / reporting adjustment & 1748 & 0.0002939 [-0.0003183, 0.0008476] & 0.01495 & 0.01466 & 1000/1000 \\
G / 2022-23 & 367 & 0.07179 [0.03786, 0.1162] & 2.702 & 37.64 & 1000/1000 \\
G / 2023-24 & 503 & 0.03046 [0.0118, 0.07462] & 2.382 & 78.2 & 1000/1000 \\
G / 2024-25 & 357 & 0.06469 [0.03191, 0.1278] & 2.082 & 32.19 & 1000/1000 \\
G / 2025-26 & 529 & 0.1383 [0.03715, 0.2252] & 0.7411 & 5.359 & 1000/1000 \\
G / reporting adjustment & 1756 & 0.000373 [-0.003559, 0.003454] & 0.08401 & 0.08364 & 1000/1000 \\
Summer 2023 & 772 & 15.55 & 94.24 & 0.4873 & -- \\
Summer 2024 & 400 & 17.12 & 45.59 & 0.6683 & -- \\
Summer 2025 & 839 & 16.25 & 92.31 & 0.4882 & -- \\
\bottomrule
\end{tabularx}
\par\smallskip
\begin{minipage}{\linewidth}\footnotesize\raggedright
Season rows report a with its interval and b, b/a as point estimates only. Reporting-adjustment rows use the a column for original-minus-adjusted a [95\% CI], b for original a, and b/a for adjusted a; the original normalizer is fixed. Summer rows use these columns for median count scale, median reporting participation (\%) and IQR(log G). The single voluntary summer is descriptive, not a causal or independent replicated intervention. RA/2022-23: 1000/1000 finite draws; RA/2023-24: 1000/1000 finite draws; RA/2024-25: 1000/1000 finite draws; RA/2025-26: 1000/1000 finite draws; G/2022-23: 1000/1000 finite draws; G/2023-24: 1000/1000 finite draws; G/2024-25: 1000/1000 finite draws; G/2025-26: 1000/1000 finite draws.
\end{minipage}
\end{table}


